\section{Introduction}

Exposure to artificial intelligence (AI) in the workforce has garnered much attention in the media due to anxiety surrounding workforce displacement. However, in crucial occupational sectors that are chronically understaffed and under-resourced, there may be more optimism about AI's potential to fulfill mission objectives with less resources. Local community service organizations that provide food, shelter, and emergency relief to people in acute need persistently operate with limited staff capacity. Recent progress in large language models (LLMs), have prompted those in the sector to wonder if these tools could relieve burden for workers in this sector. For both policymakers and frontline staff, the question is how much of the working day AI might realistically support.

A growing body of literature measures how exposed work is to artificial intelligence. Eloundou et al. (2024) rate each task in the O*NET occupational database on whether an LLM could substantially reduce completion time. Hatgis-Kessell et al. (2026) add that exposure depends not only on whether AI can perform a task but on how much time workers normally spend on the task. Massenkoff and McCrory (2026) move from theoretical to observed exposure by measuring which tasks actually see work-related AI use in real Claude queries.

These studies into occupational AI exposure are broad, economy-wide analyses. For local community relief workers, this does not address the specific, practical questions around what AI exposure means for them. Research about AI in the nonprofit sector largely comes from surveys conducted about whether AI is being used (Smith Arrillaga, Grundhoefer, and Im 2025; Shi and Azevedo 2026). In cases where these surveys report how AI is used, they describe general purposes such as documentation and research (Wong et al. 2026), without linking that data to specific workforce tasks or to the time those tasks take.

In this article, I bring together existing research into theoretical AI exposure (Eloundou et al.), time per task (Hatgis-Kessell et al.), and observed AI exposure (Massenkoff and McCrory) and apply them deeply to the three largest occupations in the sector of Community Food and Housing, and Emergency and Other Relief Services (NAICS 624200). For each of these occupations, I address three questions. First, how much of the working day is directly or partially exposed to AI? Second, does measuring AI exposure by worker time versus by task count change how exposed these occupations appear? Finally, how much worker time is already showing AI use, using observed exposure data?

This analysis shows that for the three largest jobs in the local community relief services sector, there is little direct AI exposure. For instance, among Social and Human Service Assistants, the largest occupation, about 24 minutes of the entire workday could be called directly exposed to AI. Accordingly, hopes for immediate efficiency gains from AI adoption ought to be tempered. There is greater evidence for partial AI exposure, that is, if organizations were to invest in purpose-built software to enhance LLM ability, a greater share of work could likely be offloaded. The remaining half of Social and Human Service work is unexposed (unlikely to be supported by AI) due to the work being physical, interpersonal, and/or relational in nature. This article contributes a sector-focused, time-based estimate for setting realistic expectations around what AI can return in terms of staff capacity.
